% Part: many-valued-logic
% Chapter: syntax-and-semantics
% Section: introduction

\documentclass[../../../include/open-logic-section]{subfiles}

\begin{document}

\olfileid[id]{mvl}{syn}{int}

\olsection{Pengantar}

Dalam logika klasik, kita menangani !!{formula}s yang dibangun dari
!!{propositional variable}s dengan menggunakan konektif proposisional
$\lnot$, $\land$, $\lor$, $\lif$, dan~$\liff$. Ketika mendefinisikan
semantik bagi logika klasik, kita menggunakan dua nilai kebenaran, yakni
$\True$ dan~$\False$. Kita menafsirkan !!{propositional variable}s melalui
!!a{valuation}~$\pAssign{v}$, yang menetapkan salah satu nilai kebenaran
$\True$, $\False$ kepada setiap !!{propositional variable}. Setiap
!!{valuation} kemudian menentukan suatu nilai kebenaran $\pValue{v}(!A)$
bagi setiap !!{formula}~$!A$, dan !!^a{formula} dipenuhi oleh
!!a{valuation}~$\pAssign{v}$, $\pSat{v}{!A}$, jika dan hanya jika
$\pValue{v}(!A)=\True$.

Logika bernilai banyak memperumum logika klasik bernilai dua dengan
memperbolehkan nilai kebenaran selain $\True$ dan~$\False$. Jadi, dalam
logika bernilai banyak, !!a{valuation}~$\pAssign{v}$ merupakan fungsi yang
menetapkan salah satu dari sejumlah nilai kebenaran yang mungkin kepada setiap
!!{propositional variable}~$p$. Kita pada umumnya akan menyebut himpunan nilai
kebenaran yang diperbolehkan sebagai~$V$. Logika klasik merupakan logika
bernilai banyak dengan $V=\{\True,\False\}$, dan nilai kebenaran
$\pValue{v}(!A)$ dihitung dengan menggunakan tabel kebenaran karakteristik
yang sudah dikenal bagi konektif-konektif tersebut.

Begitu kita menambahkan nilai kebenaran lain, terdapat lebih dari satu pilihan
yang wajar mengenai cara menghitung $\pValue{v}(!A)$ bagi konektif yang kita
baca sebagai ``dan'', ``atau'', ``tidak'', serta ``jika---maka''. Jadi, suatu
logika bernilai banyak tidak hanya ditentukan oleh himpunan nilai
kebenarannya, tetapi juga oleh \emph{fungsi kebenaran} yang kita pilih bagi
setiap konektif. Akan tetapi, setelah semua fungsi itu dipilih untuk suatu
logika bernilai banyak~$\Log L$, nilai kebenaran
$\pValue{v}(!A)[\Log L]$ ditentukan secara tunggal oleh valuasinya, sama
seperti dalam logika klasik. Seperti logika klasik, logika bernilai banyak
bersifat \emph{fungsional-kebenaran}.

Dengan unsur-unsur semantik ini, kita dapat mendefinisikan padanan konsep
semantik tautologi, perikutan, dan keterpenuhan. Dalam logika klasik,
!!a{formula} merupakan tautologi jika nilai kebenarannya
$\pValue{v}(!A)=\True$ bagi setiap~$\pAssign{v}$. Dalam logika bernilai
banyak, hal ini juga perlu diperumum. Pertama, himpunan nilai kebenaran~$V$
tidak disyaratkan memuat~$\True$. Sebagai contoh, beberapa logika bernilai
banyak menggunakan bilangan---misalnya semua bilangan rasional antara $0$
dan~$1$---sebagai himpunan nilai kebenarannya. Dalam hal seperti itu, $1$
biasanya mengambil peran~$\True$. Dalam logika lain, peran tersebut tidak
hanya dimainkan oleh satu nilai kebenaran, melainkan oleh beberapa nilai.
Karena itu, kita mensyaratkan agar setiap logika bernilai banyak mempunyai
suatu himpunan~$V^+$ yang terdiri atas \emph{nilai-nilai yang ditetapkan}.
Kita kemudian dapat mengatakan bahwa !!a{formula} dipenuhi oleh
!!a{valuation}~$\pAssign{v}$, $\pSat{v}{!A}[\Log L]$, jika dan hanya jika
$\pValue{v}(!A)[\Log L]\in V^+$. !!^a{formula}~$!A$ merupakan tautologi
dari logika tersebut, $\Entails[\Log L] !A$, jika dan hanya jika
$\pValue{v}(!A)[\Log L]\in V^+$ bagi setiap~$\pAssign{v}$. Terakhir, kita
mengatakan bahwa $!A$ merupakan konsekuensi logis dari suatu himpunan
!!{formula}s, $\Gamma\Entails[\Log L]!A$, jika setiap !!{valuation} yang
memenuhi semua !!{formula}s dalam~$\Gamma$ juga memenuhi~$!A$.

\end{document}
